% ========================================================================================
% CAPITOLO 5 - PROCESSORE
% ========================================================================================
\chapter{Processore}

% ----------------------------------------------------------------------------------------
\section{Esercizio 9: Processore IJVM}
% ----------------------------------------------------------------------------------------

\begin{traccia}
	A partire dall'implementazione fornita del processore operante secondo il modello IJVM,
	\begin{itemize}[nosep]
		\item si proceda all'analisi dell'architettura mediante simulazione e si approfondisca
		      lo studio del suo funzionamento per due istruzioni a scelta,
		\item si modifichi un codice operativo a scelta, documentando tutte le modifiche effettuate.
	\end{itemize}
\end{traccia}

\progetto
\todo{Schema della microarchitettura: registri, bus B e C, ALU, shifter, control store, MPC, MIR.}

\subsection{Analisi della prima istruzione}
\todo{Microprogramma commentato, tabella ciclo per ciclo, forme d'onda annotate.}

\begin{table}[H]
	\centering\small
	\caption{Esempio di tabella ciclo per ciclo per l'esecuzione di un'istruzione.}
	\label{tab:ijvm-cicli}
	\begin{tabularx}{\linewidth}{@{}c l l Y@{}}
		\toprule
		\textbf{Ciclo} & \textbf{Microistruzione} & \textbf{Registri scritti} & \textbf{Descrizione} \\ \midrule
		1 & \todo{} & & \\
		2 & & & \\
		3 & & & \\
		\bottomrule
	\end{tabularx}
\end{table}

\subsection{Analisi della seconda istruzione}
\todo{Come sopra.}

\subsection{Modifica di un codice operativo}
\todo{File modificati, righe cambiate, motivazione, programma di test prima e dopo.}
